\section{A language map for research papers}\label{appC:sec:language}
Research papers use a small vocabulary of words that are rarely defined in the paper itself: \emph{compact}, \emph{dense}, \emph{continuous}, \emph{countable}, and many others. This section is a dictionary, not a course. Each entry gives the meaning the word has in a typical paper, together with the place in this book where the idea is used or developed. An entry marked ``outside this book'' belongs to a subject you will meet next; its one-line meaning is enough to read a statement, not enough to prove something with the word.

Several phrases are already explained in the main text, and the table does not repeat them: \emph{arbitrary} (Sections~\ref{ch02:sec:quantifiers} and~\ref{ch03:sec:sets}), \emph{necessary}, \emph{sufficient}, \emph{if and only if}, \emph{converse}, \emph{contrapositive} and \emph{counterexample} (Chapter~\ref{ch:logic}), \emph{well-defined} (Section~\ref{ch06:sec:well-defined}), \emph{without loss of generality} and \emph{by symmetry} (Section~\ref{ch06:sec:wlog}), \emph{maximum}, \emph{minimum}, \emph{maximal} and \emph{minimal} (Sections~\ref{ch06:sec:order} and~\ref{ch11:sec:bounds}), \emph{bounded}, \emph{supremum} and \emph{infimum} (Section~\ref{ch11:sec:bounds}), \emph{eventually} (Section~\ref{ch11:sec:convergence}), \emph{monotone} and \emph{limits superior and inferior} (Section~\ref{ch11:sec:bounds}), and \emph{sharp} (Section~\ref{ch18:sec:sharp}).

The same is true of the basic objects of the book: function, domain, codomain, image, preimage, injection, surjection and bijection (Chapter~\ref{ch:maps}); relation, equivalence relation, equivalence class, quotient and partial order (Chapter~\ref{ch:structure}); graph, walk, path and connected graph (Chapter~\ref{ch:graphs}); vector, linear map, kernel, basis, span, inner product and norm (Chapter~\ref{ch:linear}); metric, ball, convergence, Cauchy sequence, completeness, subsequence and total boundedness (Chapters~\ref{ch:limits} and~\ref{ch:bridge}); and group and isomorphism (Chapter~\ref{ch:structure}).

\begin{longtable}{@{}>{\raggedright\arraybackslash}p{1.35in}>{\raggedright\arraybackslash}p{3.85in}@{}}
\toprule Term or phrase & What it means\\\midrule\endhead
\emph{It suffices to show} & To prove the stated claim, it is enough to prove a second statement that implies it. The second statement may be stronger, weaker, or a special case; the proof must say why it gives the original claim (Section~\ref{ch02:sec:implication}).\\
\emph{It remains to show} & The proof has established part of the claim; the named statement is the only outstanding obligation.\\
\emph{Fixed} & ``Fix $x$'' means choose $x$ once and keep it unchanged for the rest of the argument; it does not mean that $x$ is unique or that a quantity is constant (Sections~\ref{ch02:sec:quantifiers} and~\ref{ch03:sec:sets}).\\
\emph{Conversely} & Introduces the reverse implication: after proving $P\Rightarrow Q$, the word begins the proof of $Q\Rightarrow P$ (Section~\ref{ch02:sec:implication}). The phrase ``and vice versa'' says that the reverse direction also holds.\\
\emph{The following are equivalent} & A list of statements each of which implies all the others. To prove the list, it is enough to prove a cycle of implications, such as $P\Rightarrow Q\Rightarrow R\Rightarrow P$ (Section~\ref{ch03:sec:injsurj}).\\
\emph{Trivial} & Follows immediately from the definitions or from a hypothesis, without an idea. The word describes the argument, not the importance of the result. If a reader cannot see why it is trivial, the writer should supply the missing line.\\
\emph{Nontrivial, nondegenerate} & Not a trivial or exceptional case; the exact meaning depends on the subject, for example a nonzero vector, a matrix with nonzero determinant, or an interval of positive length. Outside this book.\\
\emph{Standard, classical} & A result or argument already established in the literature or earlier in the text. A careful proof names the result and checks its hypotheses; Section~\ref{ch04:sec:promises} explains why an unnamed ``standard result'' is not enough.\\
\emph{One checks, direct computation, straightforward} & A verification by substituting the definitions and simplifying. The writer should show enough intermediate steps that a reader can check the claim without guessing.\\
\emph{Well-posed, stable} & A problem is \emph{well-posed} if a solution exists, is unique, and depends continuously on the data. \emph{Stable} means that small changes in the data produce small changes in the solution. Outside this book.\\
\emph{Ansatz} & A proposed form for a solution, chosen to simplify the problem and then checked against the equations. Outside this book.\\
\emph{Remark} & A short comment that is not needed for the proof: a variant, an example, a warning, or a limitation.\\
\emph{Open problem} & A question that is not known to be true or false; a conjecture is a proposed answer to one (Section~\ref{ch01:sec:roles}).\\
\emph{Convention} & An agreed choice made for convenience, not a claim to prove; for example, $a^0=1$ for $a\ne0$ (Section~\ref{ch01:sec:letters}), or names of hypotheses begin with \code{h} in Lean (Section~\ref{ch15:sec:symbols}).\\
\emph{Abuse of notation} & A convenient shorthand that is not literally correct but causes no ambiguity, such as writing $f(x)$ for $f(\{x\})$, or using one symbol for an object and its class. Say what is meant when confusion is possible.\\
\emph{Identify} & To treat two different descriptions as the same object through a specified bijection or isomorphism; say which correspondence is used (Section~\ref{ch06:sec:structure}).\\
\emph{Canonical} & A standard choice determined by the structure of the object, not by an arbitrary decision. For example, the standard basis $e_1,\ldots,e_n$ of $\RR^n$ is canonical; a basis chosen by picking vectors at random is not (Section~\ref{ch10:sec:vectors}).\\
\emph{Up to, modulo} & ``$A$ and $B$ agree up to $X$'' means that they differ only by the thing $X$, which the discussion is ignoring. The phrase ``modulo $X$'' means the same. Say what $X$ is: up to sign, up to relabeling, up to isomorphism, modulo a constant (Section~\ref{ch06:sec:structure}). The word \emph{essentially} often signals such an ignored difference.\\
\emph{Total order, chain} & A partial order is a \emph{total order} if every two elements are comparable. In order theory, a \emph{chain} is a subset in which every two elements are comparable; Chapter~\ref{ch:graphs} uses ``chain'' for a different object, a directed path. Outside this book beyond partial orders (Section~\ref{ch06:sec:order}).\\
\emph{Closed under} & A set is \emph{closed under} an operation if applying the operation to members always gives a member of the set; for example, the integers are closed under addition but not under division. The symmetries of an object are closed under composition (Section~\ref{ch06:sec:structure}).\\
\emph{Generic} & Typical, in a sense that depends on the subject: the property holds for all objects except a small exceptional set. Outside this book.\\
\emph{Sharp, tight, optimal} & A bound is \emph{sharp} if no smaller bound holds in general, and it is \emph{attained} if equality holds for some example in the class. An inequality or a case is \emph{tight} if equality can occur; a group of students is tight when it has no spare task. An \emph{optimal} bound or object is the best possible under the stated requirements (Sections~\ref{ch09:sec:inductive} and~\ref{ch18:sec:sharp}).\\
\emph{Countable, uncountable} & A set is \emph{countable} if it is finite or its members can be listed as a sequence; otherwise it is \emph{uncountable}. The integers and the rationals are countable, while the real numbers are uncountable. Outside this book.\\
\emph{Cardinality} & The size of a set. For a finite set it is the number of members, written $|A|$ (Section~\ref{ch03:sec:sets}); for infinite sets, two sets have the same cardinality when a bijection exists between them. Outside this book beyond finite sets.\\
\emph{Power set} & The set of all subsets of a set $A$, written $2^A$ or $\mathcal P(A)$. Outside this book.\\
\emph{Almost all} & Usually ``all but finitely many''; in measure theory, ``all except a set of measure zero.'' The context decides. Outside this book.\\
\emph{a priori, a posteriori} & ``From what comes before'' and ``from what comes after'': an \emph{a priori} bound is computed from the assumptions before the computation; an \emph{a posteriori} bound uses the computed values. Chapter~\ref{ch:contraction} uses both.\\
\emph{Open, closed} & In a metric space, a set $U$ is \emph{open} if every point of $U$ has a ball around it contained in $U$. A set $C$ is \emph{closed} if its complement is open, equivalently if every convergent sequence from $C$ has its limit in $C$ (Section~\ref{ch11:sec:complete}).\\
\emph{Closure, interior, boundary} & The \emph{closure} of $A$ is the smallest closed set containing $A$; the \emph{interior} is the largest open set contained in $A$; the \emph{boundary} is the closure with the interior removed. Outside this book.\\
\emph{Limit point} & A point $x$ is a \emph{limit point} of $A$ if every ball around $x$ contains a point of $A$ other than $x$; equivalently, some sequence from $A\setminus\{x\}$ converges to $x$. The closure of $A$ is $A$ together with its limit points. Outside this book.\\
\emph{Dense} & A set $D$ is dense in a space $X$ if every ball in $X$ contains a point of $D$, equivalently if the closure of $D$ is $X$. The rationals are dense in the real line. Outside this book.\\
\emph{Compact} & In a metric space, \emph{compact} means that every sequence has a convergent subsequence; equivalently, every open cover (a collection of open sets whose union is the space) has a finite subcover. The theorem of Chapter~\ref{ch:bridge} proves that a complete, totally bounded metric space has the first property (Section~\ref{ch17:sec:proof}).\\
\emph{Continuous} & A map $f:X\to Y$ between metric spaces is \emph{continuous} if $x_n\to x$ always implies $f(x_n)\to f(x)$; equivalently, the preimage of every open set is open. Outside this book.\\
\emph{Connected} & A space is \emph{connected} if it cannot be split into two disjoint nonempty open sets. Outside this book.\\
\emph{Neighborhood} & A \emph{neighborhood} of a point is a set containing a ball around that point. The \emph{out-neighborhood} of a vertex in Chapter~\ref{ch:graphs} is a different, graph-specific use of the word (Section~\ref{ch07:sec:seymour}).\\
\emph{Totally bounded} & For every radius $r>0$, finitely many balls of radius $r$ cover the space (Section~\ref{ch17:sec:balls}).\\
\emph{Pointwise, uniform} & A property holds \emph{pointwise} if it holds at each object separately, with a choice that may depend on the object; it holds \emph{uniformly} if one choice works for all objects at once. Pointwise is used in Section~\ref{ch13:sec:expectation}; uniform convergence is outside this book.\\
\emph{$O$, $o$, $\Theta$, $\Omega$, $\sim$} & For sequences with $g(n)\ne0$ eventually, $f(n)=O(g(n))$ means $|f(n)|\le C|g(n)|$ for some constant $C>0$ and all sufficiently large $n$; $f(n)=o(g(n))$ means $f(n)/g(n)\to0$; $f(n)=\Omega(g(n))$ means $g(n)=O(f(n))$; $f(n)=\Theta(g(n))$ means both $O$ and $\Omega$; and $f(n)\sim g(n)$ means $f(n)/g(n)\to1$. Outside this book, which states its bounds explicitly.\\
\emph{a.e., a.s.} & A property holds \emph{almost everywhere} if it fails only on a set of measure zero, that is, a set that can be covered by intervals of arbitrarily small total length. In probability, \emph{almost surely} (a.s.) means with probability $1$. Measure theory is outside this book.\\
\emph{i.i.d.} & Independent and identically distributed: the random variables have the same distribution, and events determined by disjoint groups of them are independent. Outside this book.\\
\emph{iff, s.t., w.r.t., w.l.o.g.} & Shorthands for ``if and only if,'' ``such that,'' ``with respect to,'' and ``without loss of generality.'' This book writes them in full (Sections~\ref{ch02:sec:implication} and~\ref{ch06:sec:wlog}).\\
\emph{i.e., e.g., cf., resp.} & ``That is,'' ``for example,'' ``compare,'' and ``respectively.'' This book spells these out.\\
\emph{QED, $\square$} & The end of a proof. The square is used in this book (Section~\ref{ch01:sec:roles}).\\
\bottomrule
\end{longtable}

A dictionary entry is enough to read a sentence containing the word, but not enough to use the word in a proof. When a paper makes one of these words part of a hypothesis, look up its definition in that subject, and check the examples and boundary cases that the definition allows.

\subsection*{Standard argument patterns}
Some sentences in papers name an argument rather than a result. Four patterns are common enough to record. Each is a recipe for a proof, not a substitute for the definitions it uses.

\emph{By density and continuity.} Prove the statement on a dense subset, then extend it to the whole space because both sides are continuous. If two continuous functions agree on a dense set, they agree everywhere.

\emph{By compactness.} In a compact setting, extract a convergent subsequence from a sequence, or a finite subcover from an open cover, and pass to the limit. In a compact metric space every sequence has a convergent subsequence.

\emph{By analyticity.} If two analytic functions agree on a set with an accumulation point in a connected domain, they agree on the whole domain. This is why analytic continuation is unique.

\emph{By a finite-net argument.} Cover a compact set by finitely many small balls, check an estimate on each ball, and combine the finitely many bounds into one uniform bound.

\subsection*{Where this map stops}
This map stops at the border of general language. Papers in analysis, spectral theory and numerical analysis use a further layer of domain vocabulary: operator, adjoint, self-adjoint, resolvent, spectrum, Fredholm, Sobolev, weak convergence, Galerkin, and many more. Those words are not general grammar; they belong to those subjects. The grammar in this book is enough to ask for their definitions, to see where they enter a statement, and to check the examples and boundary cases; it is not a substitute for learning the definitions from a course or an AI tutor.
